\documentclass[11pt]{amsart}

\usepackage[T1]{fontenc}
\usepackage[utf8]{inputenc}
\usepackage{amsmath,amssymb,amsthm}
\usepackage[margin=1in]{geometry}
\usepackage{array}
\usepackage{booktabs}
\usepackage{enumitem}
\usepackage{longtable}
\usepackage{hyperref}

\newcommand{\LL}{\Lambda}
\newcommand{\OO}{\mathbb{O}}
\newcommand{\ZZ}{\mathbb{Z}}
\newcommand{\QQ}{\mathbb{Q}}
\newcommand{\FF}{\mathbb{F}}
\newcommand{\RR}{\mathbb{R}}
\newcommand{\CC}{\mathbb{C}}
\newcommand{\Nrm}{N}
\newcommand{\ip}[2]{\langle #1, #2 \rangle}
\DeclareMathOperator{\Aut}{Aut}
\DeclareMathOperator{\Der}{Der}
\DeclareMathOperator{\Ann}{Ann}
\DeclareMathOperator{\Stab}{Stab}
\DeclareMathOperator{\End}{End}
\DeclareMathOperator{\Co}{Co}
\DeclareMathOperator{\Tr}{tr}
\DeclareMathOperator{\Min}{Min}
\DeclareMathOperator{\Rea}{Re}

\newtheorem{theorem}{Theorem}
\newtheorem{proposition}[theorem]{Proposition}
\newtheorem{lemma}[theorem]{Lemma}
\newtheorem{corollary}[theorem]{Corollary}
\theoremstyle{definition}
\newtheorem{definition}[theorem]{Definition}
\theoremstyle{remark}
\newtheorem{remark}[theorem]{Remark}

\title[The automorphism group of the triple-octonion product]%
{The automorphism group of the\\
$\ZZ_3$-symmetric triple-octonion product\\
\large Companion calculation to
\emph{An order on the Leech lattice from a $\ZZ_3$-symmetric
triple-octonion product}}

\author{Claude Opus 4.8}
\address{Anthropic}
\date{12 July 2026}

\thanks{This note was computed at the direction of Jens K\"oplinger and is
companion material to the main paper \emph{An order on the Leech lattice
from a $\ZZ_3$-symmetric triple-octonion product} (v6). It is not part of
that paper and is not in its referee scope; the author of the main paper
will decide how, and how much of, this material is worked into it.
Every nontrivial statement below is either proved here or reproduced by a
named script in the project repository. Two independent computational
passes were run: a first pass that produced the result, and a second pass
that rebuilt the lattice, the product and the generators from the paper's
definitions without reusing any code from the first. Where the two passes
disagreed, the disagreement is recorded in \S\ref{sec:corrections}.}

\setlength{\emergencystretch}{3em}

\begin{document}
\maketitle

\begin{abstract}
The main paper leaves open the identification of $\Aut(\LL, +, \star)$,
the group of $\ZZ$-linear bijections of the Leech lattice that preserve the
$\ZZ_3$-symmetric $\sigma$-twisted triple-octonion product~$\star$.
That group has order~$36$ and is isomorphic to $C_6 \times S_3$: every
element is a blockwise octonion automorphism composed with a permutation
of the three octonion blocks, the blockwise part running over a cyclic
group of order~$6$ of signed permutations of the octonion basis.
It is contained in the Conway group~$\Co_0$, and this containment is a
theorem rather than an observation: the ambient algebraic group
$\Aut(\RR^{24}, \star)$ is the compact group $G_2 \times G_2 \times S_3$
of dimension~$28$, and compactness forces every $\star$-automorphism to be
a Euclidean isometry.
\end{abstract}

% ===================================================================
\section{Setting}\label{sec:setting}
% ===================================================================

The objects are those of the main paper, restated here only far enough to
make this note self-contained.

$\OO$ is the real octonion algebra on $\RR^8$ with basis
$e_0, \ldots, e_7$, unit~$e_0$, and the standard Fano triples
$(1,2,4)$, $(2,3,5)$, $(3,4,6)$, $(4,5,7)$, $(5,6,1)$, $(6,7,2)$,
$(7,1,3)$: for each cyclic triple $(a,b,c)$ one has $e_a e_b = +e_c$, for
each anticyclic one $e_a e_b = -e_c$, and $e_i^2 = -e_0$ for
$i = 1, \ldots, 7$. Write $\Nrm$ for the norm form and $\Rea(x)$ for the
$e_0$-component of~$x$.

$\sigma = (1\;2)$ is the linear involution of $\RR^8$ fixing $e_0$ and
interchanging $e_1$ and~$e_2$. The $\sigma$-twisted octonion product is
\[
  x \cdot_s y \;:=\; \sigma\bigl(\sigma(x)\,\sigma(y)\bigr),
\]
which is again an octonion product on $\RR^8$, and $\sigma$ is by
construction an algebra isomorphism
$(\OO, \cdot) \to (\OO, \cdot_s)$.

$L = D_8^+$ is the $E_8$ lattice realized in octonion coordinates,
$s = \tfrac12(-e_0 + e_1 + \cdots + e_7)$,
$\bar s = \tfrac12(-e_0 - e_1 - \cdots - e_7)$, and
$Ls = \{\ell \cdot s : \ell \in L\}$,
$L\bar s = \{\ell \cdot \bar s : \ell \in L\}$, both formed with the
\emph{untwisted} octonion product. Wilson's description of the Leech
lattice inside $L^3$ is
\begin{equation}\label{eq:wilson}
  \LL \;=\; \bigl\{ (x,y,z) \in L^3 \;:\;
    x + y,\; x + z,\; y + z \in L\bar s; \;\;
    x + y + z \in Ls \bigr\}.
\end{equation}

On $\RR^{24} = \OO_1 \oplus \OO_2 \oplus \OO_3$ the product of the main
paper is
\begin{equation}\label{eq:star}
  (x,y,z) \star (x',y',z') \;=\;
  \bigl(\,
    x \cdot_s x' + z \cdot_s y' + y \cdot_s z',\;\;
    y \cdot_s y' + x \cdot_s z' + z \cdot_s x',\;\;
    z \cdot_s z' + y \cdot_s x' + x \cdot_s y'
  \,\bigr),
\end{equation}
and the main paper's Theorem~1.1 is $\LL \star \LL \subseteq \LL$.

Throughout, blocks are indexed by $\ZZ/3 = \{0,1,2\}$ and a vector of
$\RR^{24}$ is written $u = (u_0, u_1, u_2)$ with $u_\alpha \in \OO$. The
object under study is
\[
  \Aut(\LL, +, \star) \;:=\;
  \bigl\{\, g : \LL \to \LL \;\text{a $\ZZ$-linear bijection with}\;
    g(u \star v) = g(u) \star g(v) \;\text{for all}\; u,v \in \LL \,\bigr\}.
\]
The ambient group here is $\mathrm{GL}_{24}(\ZZ)$ and not $\Co_0$: that
every such~$g$ is in fact an isometry is one of the things settled below
(\S\ref{sec:co0}), not something assumed.

Two conventions should be flagged. First, $\star$ is bilinear and
$\LL$-membership is $\ZZ$-linear, so a check carried out on all
$24 \times 24 = 576$ pairs of basis vectors is a \emph{proof} for all of
$\RR^{24}$, not a sample; every ``exhaustive basis check'' below is of
this kind. Second, $\LL$ as realized by~\eqref{eq:wilson} has minimal norm
$8$, Gram determinant $2^{24}$ and index $4096$ in~$L^3$: it is a copy of
$\sqrt2\,\LL_{\mathrm{Leech}}$ rather than the unimodular Leech lattice.
Nothing below depends on the scale ($\Co_0$ is scale invariant), but no
claim of unimodularity should be attached to this realization.

% ===================================================================
\section{The routing identity and the ideal decomposition}\label{sec:ideals}
% ===================================================================

The nine octonion blocks in~\eqref{eq:star} are routed by a single
congruence.

\begin{lemma}[Routing]\label{lem:routing}
For $u, v \in \RR^{24}$ and $\gamma \in \ZZ/3$,
\[
  (u \star v)_\gamma \;=\;
  \sum_{\substack{\alpha, \beta \in \ZZ/3 \\ \alpha + \beta \equiv 2\gamma}}
    u_\alpha \cdot_s v_\beta .
\]
\end{lemma}

\begin{proof}
Read off~\eqref{eq:star} with $(u_0,u_1,u_2) = (x,y,z)$. The first block
collects the index pairs $(0,0), (2,1), (1,2)$, all with
$\alpha + \beta \equiv 0 = 2\cdot 0$; the second collects
$(1,1), (0,2), (2,0)$, all with $\alpha + \beta \equiv 2 = 2 \cdot 1$; the
third collects $(2,2), (1,0), (0,1)$, all with
$\alpha + \beta \equiv 1 \equiv 2 \cdot 2$.
\end{proof}

Since $2$ is invertible modulo~$3$, each of the nine pairs $(\alpha,\beta)$
occurs in exactly one block, namely $\gamma \equiv 2(\alpha+\beta)$. Two
consequences are immediate.

\begin{proposition}[Ideal decomposition]\label{prop:ideals}
Let
\[
  D \;:=\; \{(a,a,a) : a \in \OO\}, \qquad
  T \;:=\; \{(p,q,r) \in \OO^3 : p + q + r = 0\},
\]
and let $\Sigma(u) := u_0 + u_1 + u_2$. Then
\begin{enumerate}[label=(\arabic*)]
\item $\Sigma(u \star v) = \Sigma(u) \cdot_s \Sigma(v)$, so $\Sigma$ is a
  surjective algebra homomorphism $(\RR^{24}, \star) \to (\OO, \cdot_s)$;
\item $(a,a,a) \star v = \bigl(a \cdot_s \Sigma(v)\bigr)^{\mathrm{diag}}$
  and $v \star (a,a,a) = \bigl(\Sigma(v) \cdot_s a\bigr)^{\mathrm{diag}}$,
  where $w^{\mathrm{diag}} := (w,w,w)$;
\item $D$ and $T = \ker \Sigma$ are two-sided ideals,
  $D \star T = T \star D = 0$, and $\RR^{24} = D \oplus T$;
\item $\varphi(a) := \tfrac13 (a,a,a)$ is an algebra isomorphism
  $(\OO, \cdot_s) \to (D, \star)$; equivalently
  $(a,a,a) \star (b,b,b) = 3\,(a \cdot_s b,\, a \cdot_s b,\, a \cdot_s b)$;
\item $D \perp T$ for the standard Euclidean form of $\RR^{24}$.
\end{enumerate}
\end{proposition}

\begin{proof}
(1) Summing Lemma~\ref{lem:routing} over $\gamma$ collects each pair
$(\alpha,\beta)$ exactly once, giving
$\sum_{\alpha,\beta} u_\alpha \cdot_s v_\beta
 = \Sigma(u) \cdot_s \Sigma(v)$ by bilinearity. Surjectivity is clear from
$\Sigma(a,0,0) = a$.
(2) With $u = (a,a,a)$ the inner sum of Lemma~\ref{lem:routing} is
$\sum_{\beta} a \cdot_s v_\beta = a \cdot_s \Sigma(v)$ for every~$\gamma$,
independent of~$\gamma$; the right-multiplication case is the same
computation.
(3) $T$ is a two-sided ideal because it is the kernel of the homomorphism
$\Sigma$; $D$ is a two-sided ideal by~(2). If $v \in T$ then $\Sigma(v)=0$,
so~(2) gives $D \star T = T \star D = 0$. Finally
$\Sigma|_D : (a,a,a) \mapsto 3a$ is bijective onto~$\OO$ and $T = \ker\Sigma$,
so $\RR^{24} = D \oplus T$.
(4) Immediate from~(2) with $v = (b,b,b)$.
(5) $\ip{(a,a,a)}{(p,q,r)} = \ip{a}{p+q+r} = 0$.
\end{proof}

All five items were additionally checked by exhaustive basis computation in
exact rational arithmetic, and re-checked against an independently coded
nine-term expansion of~$\star$
(\texttt{python\_project/src/verify\_ideal\_decomposition.py}).

\begin{remark}
The same script verifies that
$\sigma^{\oplus 3} : (x,y,z) \mapsto (\sigma x, \sigma y, \sigma z)$ is an
algebra isomorphism from the $\sigma$-twisted triple product onto the
untwisted one. The twist is invisible to the algebra: only the lattice
$\LL$ sees it. Consequently everything in
\S\ref{sec:ideals}--\S\ref{sec:ambient} could be stated without~$\sigma$;
$\sigma$ re-enters in \S\ref{sec:lattice}, where the lattice conditions do.
\end{remark}

\subsection*{Every $\star$-automorphism preserves $D$ and $T$}

Because $D \star T = T \star D = 0$, the two projections
$\pi_D, \pi_T$ associated with $\RR^{24} = D \oplus T$ are algebra
homomorphisms: for $u = u_D + u_T$ and $v = v_D + v_T$,
$u \star v = u_D \star v_D + u_T \star v_T$ with the first term in~$D$ and
the second in~$T$, so
$\pi_T(u \star v) = \pi_T(u) \star \pi_T(v)$ and likewise for~$\pi_D$.

\begin{lemma}\label{lem:simple}
$(D, \star)$ and $(T, \star)$ are simple algebras.
\end{lemma}

\begin{proof}[Proof and computation]
$D \cong (\OO, \cdot_s)$ by Proposition~\ref{prop:ideals}(4), and the real
octonions form a division algebra: if $I \subseteq \OO$ is a nonzero ideal
and $0 \neq a \in I$, then right multiplication $R_a$ is invertible and
$R_a(\OO) \subseteq I$, so $I = \OO$. This is a proof.

For $T$ the statement is established by exact linear algebra: the
multiplication algebra of~$T$, that is the unital associative subalgebra of
$\End(T)$ generated by all left and right $\star$-multiplications by
elements of~$T$, has dimension~$256 = \dim \End(T)$, so it is all of
$\End(T) \cong M_{16}(\RR)$. Any ideal of~$T$ is invariant under the
multiplication algebra, hence under all of $\End(T)$, hence is $0$ or~$T$.
The dimension is a rank of an explicit rational matrix; it was computed
exactly, with a full-rank certificate modulo a prime as a rigorous lower
bound. (Second pass, independent rebuild; the same computation gives
dimension~$64 = \dim\End(D)$ for~$D$, reproving the first paragraph.)
\end{proof}

\begin{proposition}\label{prop:preserve}
Every $g \in \Aut(\RR^{24}, \star)$ satisfies $g(D) = D$ and $g(T) = T$, and
consequently
\[
  \Aut(\RR^{24}, \star) \;=\; \Aut(D, \star) \times \Aut(T, \star),
\]
the two restrictions being independently choosable.
\end{proposition}

\begin{proof}
$g(D)$ is an ideal of $\RR^{24}$ of dimension~$8$. Applying the surjective
algebra homomorphism~$\pi_T$, the image $\pi_T(g(D))$ is an ideal of~$T$ of
dimension at most~$8 < 16$. By Lemma~\ref{lem:simple} it is therefore~$0$,
so $g(D) \subseteq \ker \pi_T = D$, and $g(D) = D$ by bijectivity.

Next, $T = \Ann(D) := \{u : u \star D = D \star u = 0\}$. Indeed
$T \subseteq \Ann(D)$ by Proposition~\ref{prop:ideals}(3); conversely, write
$u = d + t$ with $d \in D$, $t \in T$, and note that
$u \star (a,a,a) = d \star (a,a,a)$ for every~$a$, so $u \in \Ann(D)$ forces
$d$ to annihilate the division algebra $(D,\star) \cong (\OO, \cdot_s)$,
whence $d = 0$. (Confirmed as an exact nullspace computation:
$\dim \Ann(D) = 16$ and $\dim \Ann(T) = 8$.) Since $g(D) = D$ is now known,
$g(T) = g(\Ann(D)) = \Ann(g(D)) = \Ann(D) = T$.

Finally, since $D \star T = T \star D = 0$, any pair of algebra automorphisms
of $D$ and of~$T$ assembles to an automorphism of~$\RR^{24}$, and by the
above every automorphism arises this way.
\end{proof}

\begin{remark}[A correction, recorded]\label{rem:circular}
The first computational pass justified Proposition~\ref{prop:preserve} by
``$\Ann(T) = D$ and $\Ann(D) = T$; annihilators are intrinsic, so every
automorphism preserves $D$ and~$T$''. As written that is circular:
$g(\Ann(D)) = \Ann(g(D))$ yields $g(T) = T$ only \emph{after} $g(D) = D$ is
known, and the annihilator relations alone do not establish that~$D$ is
intrinsic. The simplicity argument above is the non-circular replacement and
is what the second pass verified. The conclusion is unaffected. The
docstring of \texttt{python\_project/src/verify\_star\_algebra\_structure.py}
still carries the circular wording and should be corrected.
\end{remark}

% ===================================================================
\section{The $C_3$-Fourier normal form}\label{sec:fourier}
% ===================================================================

The routing identity diagonalizes under the discrete Fourier transform of
the block index. Let $\zeta$ be a primitive cube root of unity and, for
$u \in \RR^{24} \otimes \CC$, put
\[
  \hat u_j \;:=\; \sum_{\alpha \in \ZZ/3} \zeta^{\,j\alpha}\, u_\alpha
  \;\in\; \OO_\CC, \qquad j \in \ZZ/3 .
\]

\begin{proposition}\label{prop:fourier}
For all $u, v$,
\[
  \widehat{(u \star v)}_k \;=\; \hat u_{2k} \cdot_s \hat v_{2k},
  \qquad k \in \ZZ/3 .
\]
\end{proposition}

\begin{proof}
By Lemma~\ref{lem:routing}, the block $\gamma$ receiving the pair
$(\alpha,\beta)$ is $\gamma \equiv 2(\alpha+\beta)$. Hence
\[
  \widehat{(u \star v)}_k
  = \sum_{\gamma} \zeta^{\,k\gamma} (u \star v)_\gamma
  = \sum_{\alpha, \beta} \zeta^{\,2k(\alpha+\beta)}\, u_\alpha \cdot_s v_\beta
  = \Bigl(\sum_\alpha \zeta^{\,2k\alpha} u_\alpha\Bigr) \cdot_s
    \Bigl(\sum_\beta \zeta^{\,2k\beta} v_\beta\Bigr)
  = \hat u_{2k} \cdot_s \hat v_{2k}. \qedhere
\]
\end{proof}

The identity was also verified exactly on all $576$ basis pairs over
$\QQ(\zeta)$, elements being carried as pairs $p + q\zeta$ with
$\zeta^2 = -1 - \zeta$
(\texttt{python\_project/src/verify\_star\_algebra\_structure.py}); the
second pass verified the equivalent statement in its own coordinates.

The case $k = 0$ is Proposition~\ref{prop:ideals}(1). The cases $k = 1, 2$
say that on $T \otimes \CC$, with coordinates $(\hat u_1, \hat u_2) \in
\OO_\CC \oplus \OO_\CC$, the product is the \emph{swap-twisted} algebra
\begin{equation}\label{eq:swap}
  S : \qquad (a_1, a_2) \cdot (b_1, b_2) \;=\;
    (a_2 \cdot_s b_2,\; a_1 \cdot_s b_1)
  \qquad \text{on } \OO_\CC \oplus \OO_\CC ,
\end{equation}
and that the real structure on $S$ is
$(a_1, a_2) \mapsto (\overline{a_2}, \overline{a_1})$, since
$\hat u_2 = \overline{\hat u_1}$ for real~$u$.

This is the structural engine of the calculation: the $16$-dimensional
piece~$T$, which is where all the interesting automorphisms live, is a real
form of a pair of complexified octonion algebras glued by a swap.

% ===================================================================
\section{The automorphism group of the ambient algebra}\label{sec:ambient}
% ===================================================================

For $A, B \in \Aut(\OO, \cdot_s)$ define $h_{A,B} := P \otimes A + Q \otimes B$,
where $P$ is the orthogonal projection of $\RR^{24}$ onto~$D$ and
$Q = I - P$ the projection onto~$T$; concretely
\begin{equation}\label{eq:hAB}
  h_{A,B}(u)_\alpha \;=\; B\,u_\alpha
    \;+\; (A - B)\,\tfrac{1}{3}\,(u_0 + u_1 + u_2).
\end{equation}
For $\pi \in S_3$ let $\rho_\pi(u)_\alpha := u_{\pi^{-1}(\alpha)}$ be the
corresponding permutation of the three blocks.

\begin{theorem}\label{thm:ambient}
\[
  \Aut(\RR^{24}, \star)
  \;=\; \{\, h_{A,B}\, \rho_\pi \;:\;
      A, B \in \Aut(\OO, \cdot_s),\; \pi \in S_3 \,\}
  \;\cong\; G_2 \times G_2 \times S_3 ,
\]
a \emph{direct} product, where $G_2$ denotes the compact real form. The
group is compact of dimension~$28$; $h_{A,B}$ and $\rho_\pi$ commute.
\end{theorem}

The proof has three parts.

\medskip
\noindent\textbf{(a) The $D$ factor.}
$\Aut(D, \star) = \Aut(\OO, \cdot_s)$ by
Proposition~\ref{prop:ideals}(4), and $\Aut(\OO, \cdot_s)$ is the compact
real form of~$G_2$ (classical; see also \S\ref{sec:co0}). The rescaling by~$3$
does not affect the automorphism group.

\medskip
\noindent\textbf{(b) The identity component of the $T$ factor.}
Exact nullspace computations give $\dim \Der(\OO, \cdot_s) = 14$ and
$\dim \Der(\RR^{24}, \star) = 28$: $28$ explicit derivations of the form
$P \otimes X + Q \otimes Y$ with $X, Y \in \Der(\OO, \cdot_s)$ are exhibited
and shown to be linearly independent over~$\QQ$, giving $\dim \geq 28$, and
a rank bound modulo a prime (full rank over $\FF_p$ is a rigorous lower
bound for the rank over~$\QQ$, hence the mod-$p$ nullity of the derivation
system is a rigorous upper bound for its rank over~$\QQ$) gives
$\dim \leq 28$. Hence
$\Der(\RR^{24}, \star) = \mathfrak{g}_2 \oplus \mathfrak{g}_2$; the second
pass additionally computed the Killing form of $\Der(\OO, \cdot_s)$ and found
it negative definite, so both summands are of compact type. In particular
$\Aut(T, \star)^{\,0}$ is the blockwise compact~$G_2$.
(\texttt{verify\_star\_algebra\_structure.py}; second pass independently.)

\medskip
\noindent\textbf{(c) The component group of the $T$ factor.}
Two independent derivations agree.

\emph{Route~1 (over $\CC$, via~\eqref{eq:swap}).}
In $S$, for $w = (a_1, a_2)$ the left multiplication $L_w$ has rank~$4$
exactly when one component vanishes and the other is a nonzero null vector
of~$\OO_\CC$. An automorphism preserves that rank condition, so it permutes
the two irreducible punctured null cones $N_1, N_2$ and hence their
spans~$V_1, V_2$. The component-preserving automorphisms solve out exactly to
$(\phi_1, \phi_2) = (c^2\theta,\, c\theta)$ with
$\theta \in \Aut(\OO_\CC)$ and $c^3 = 1$; the swap
$\tau(a_1,a_2) = (a_2,a_1)$ is an automorphism, fixes~$\theta$ and inverts~$c$.
Hence $\Aut(S) = G_2(\CC) \times S_3$. Passing to real points, the real
structure $(a_1,a_2) \mapsto (\overline{a_2}, \overline{a_1})$ forces
$\theta$ to be fixed by conjugation, that is $\theta$ lies in the compact
$G_2$, while $c$ is unconstrained among the cube roots of unity. So
$\Aut(T, \star)(\RR) = G_2 \times S_3$, the $G_2$ acting blockwise and the
$S_3$ being the block permutations restricted to~$T$.

\emph{Route~2 (finite and machine-checkable).}
$\Aut(G_2) = \mathrm{Inn}(G_2)$ and $Z(G_2) = 1$, so the component group of
$\Aut(T,\star)$ is the centralizer of the identity component
$\Aut(T,\star)^{\,0}$ inside $\Aut(T,\star)$. The commutant of the blockwise
$\mathfrak{g}_2$ inside $\End(T)$ was computed exactly and is
$8$-dimensional ($\cong M_2 \oplus M_2$); imposing the automorphism condition
on it yields $32$ distinct polynomial equations in $8$ unknowns, solved
exactly, with exactly $6$ invertible solutions. They are closed under
composition and have orders $1, 2, 2, 2, 3, 3$, hence form~$S_3$. There are
no further components.

Route~2 was carried out in the second, independent pass and does not share a
step with Route~1. Route~1 is the derivation recorded in
\texttt{verify\_star\_algebra\_structure.py}. The two agree.

\medskip
\noindent\textbf{Assembly.}
By Proposition~\ref{prop:preserve},
$\Aut(\RR^{24},\star) = \Aut(D,\star) \times \Aut(T,\star)
= G_2 \times (G_2 \times S_3)$. Writing an element as $h_{A,B}\rho_\pi$ uses
$A$ for the action on~$D$ and $B$ for the blockwise action on~$T$; the block
permutations act trivially on~$D$ and commute with blockwise maps, so the
product is direct and $h_{A,B}\rho_\pi = \rho_\pi h_{A,B}$. That every
$h_{A,B}$ and every $\rho_\pi$ \emph{is} a $\star$-automorphism was checked
exhaustively on all $576$ basis pairs. $\square$

\begin{remark}[What deserves a human reading]\label{rem:humanread}
Everything in Theorem~\ref{thm:ambient} is either an exact finite computation
or elementary, with one exception: the classification of $\Aut(S)$ over~$\CC$
in Route~1, specifically the rank-of-$L_w$ / null-cone step. It is the only
genuinely non-mechanical argument in this note. The second pass deliberately
did \emph{not} reproduce it, taking Route~2 instead, and reached the same
group. If Route~1 is printed anywhere, it should be re-read by hand; Route~2
is the safer thing to publish.
\end{remark}

% ===================================================================
\section{The lattice stabilizer}\label{sec:lattice}
% ===================================================================

Since $\LL$ spans $\RR^{24}$, a $\ZZ$-linear bijection of~$\LL$ extends
uniquely to $\RR^{24}$, and $\star$-multiplicativity on~$\LL$ extends by
bilinearity to all of~$\RR^{24}$. Hence
\[
  \Aut(\LL, +, \star)
  \;=\; \{\, g \in \Aut(\RR^{24}, \star) : g(\LL) = \LL \,\}
  \;=\; \{\, h_{A,B}\, \rho_\pi \;:\; h_{A,B}(\LL) = \LL,\ \pi \in S_3 \,\},
\]
the last equality because every $\rho_\pi$ preserves $\LL$: Wilson's three
conditions in~\eqref{eq:wilson} are symmetric in $x, y, z$. So the whole
problem is the determination of
\[
  K \;:=\; \bigl\{ (A,B) \in \Aut(\OO,\cdot_s)^2 \;:\;
    h_{A,B}(\LL) = \LL \bigr\},
  \qquad
  \Aut(\LL,+,\star) \;\cong\; K \times S_3 .
\]

\subsection{The two lattice slices}

\begin{lemma}\label{lem:slices}
$\LL \cap T = \{(p,q,r) \in (L\bar s)^3 : p + q + r = 0\}$, and
$\LL \cap D = \{(a,a,a) : a \in Ls\}$.
\end{lemma}

\begin{proof}[Proof of the first, computation of the second]
If $(p,q,r) \in \LL$ with $p+q+r = 0$, then $p + q = -r \in L\bar s$ and
likewise for the other two pairwise sums, and $L\bar s \subseteq L$, so
$p, q, r \in L\bar s$; the condition $p+q+r = 0 \in Ls$ is automatic. The
converse is the same reading. This is a proof.

For the second slice, $(a,a,a) \in \LL$ says $a \in L$, $2a \in L\bar s$ and
$3a \in Ls$. All three conditions are constant on the classes of $L/2L$
(because $2L \subseteq L\bar s$ and $2L \subseteq Ls$, both verified), so
they may be tested on the $256$ classes; exactly $16$ survive, and they are
precisely the classes of $Ls$, which has index $16$ in~$L$. The inclusion
$\supseteq$ is immediate ($a \in Ls$ gives $3a \in Ls$ and $2a \in 2L
\subseteq L\bar s$); the reverse inclusion is the computation.
(\texttt{python\_project/src/verify\_aut\_lambda\_star.py}, reproduced in the
second pass from an independently built $\ZZ$-basis of~$\LL$.)
\end{proof}

\begin{corollary}[Finiteness]\label{cor:finite}
If $h_{A,B}\rho_\pi \in \Aut(\LL,+,\star)$ then $A(Ls) = Ls$ and
$B(L\bar s) = L\bar s$. Both stabilizers are finite, so
$\Aut(\LL,+,\star)$ is finite.
\end{corollary}

\begin{proof}
$h_{A,B}$ acts on $D$ as $(a,a,a) \mapsto (Aa,Aa,Aa)$ and preserves $D$ and
$T$; block permutations act trivially on~$D$. By Lemma~\ref{lem:slices},
preserving $\LL \cap D$ forces $A(Ls) = Ls$. On $T$, $h_{A,B}$ acts
blockwise by~$B$; every $p \in L\bar s$ occurs as the first component of
$(p,-p,0) \in \LL \cap T$, so preserving $\LL \cap T$ forces
$B(L\bar s) = L\bar s$. A lattice stabilizer inside the orthogonal group is
finite. (Alternatively: $\Aut(\RR^{24},\star)$ is compact by
Theorem~\ref{thm:ambient} and $\Aut(\LL,+,\star)$ is a discrete subgroup of
it.)
\end{proof}

\subsection{The three octonion-automorphism stabilizers}

Corollary~\ref{cor:finite} reduces the search to the two finite groups
\[
  \Stab(Ls) := \Aut(\OO,\cdot_s) \cap \{A : A(Ls) = Ls\},
  \qquad
  \Stab(L\bar s) := \Aut(\OO,\cdot_s) \cap \{B : B(L\bar s) = L\bar s\},
\]
and it is convenient to have $\Stab(L)$ as well. All three were enumerated
completely, in exact arithmetic.

\begin{center}
\begin{tabular}{lrl}
\toprule
Group & Order & Structure (GAP) \\
\midrule
$\Stab(L)$        & $1344$  & $(C_2 \times C_2 \times C_2).\mathrm{PSL}(3,2)$ \\
$\Stab(Ls)$       & $168$   & $(C_2 \times C_2 \times C_2){:}(C_7{:}C_3)$ \\
$\Stab(L\bar s)$  & $12096$ & $\mathrm{PSU}(3,3){:}C_2 = U_3(3).2 = G_2(2)$ \\
\midrule
$\Stab(L) \cap \Stab(Ls) \cap \Stab(L\bar s)$ & $6$ & $C_6$ \\
\bottomrule
\end{tabular}
\end{center}

The identifications are GAP's, from
$8 \times 8$ generators exported by the Python enumeration
(\texttt{octonion\_stabilisers.g}); the $12096$ was additionally
confirmed isomorphic to the ATLAS group $U_3(3).2$.

Completeness of the enumerations rests on the following. An automorphism
of~$\OO$ fixes~$e_0$, is orthogonal (\S\ref{sec:co0}), permutes the minimal
vectors of any lattice it stabilizes, and is determined by the images of
\emph{three} elements that generate $\OO$ as an algebra. Three, not two: by
Artin's theorem any two octonions generate an associative, hence at most
quaternionic, subalgebra, so a search over the images of two generators
undercounts the constraints and returns a much larger set. A backtracking
search over ordered generating triples drawn from the minimal-vector shell,
with Gram-matrix filters, is therefore exhaustive.

The enumerations were checked in three further ways.
(i)~Re-running the backtracking search from a different first generator, which
produces a different generating triple, different Gram filters and a different
candidate list, returns the identical \emph{set} of matrices in all three
cases (\texttt{python\_project/src/verify\_aut\_octonion\_crosscheck.py}).
(ii)~For $\Stab(L)$ an assumption-free brute force over all
$2^7 \cdot 7! = 645{,}120$ signed permutations of $e_1, \ldots, e_7$ returns
the same $1344$ matrices. This is legitimate because every element of
$\Stab(L)$ \emph{is} such a signed permutation: $e_0 \pm e_j$ are roots of~$L$
and are the only roots of $L$ of the form $e_0 + (\text{imaginary unit})$, so
$A(e_0 \pm e_j) = e_0 \pm A(e_j)$ forces $A(e_j) = \pm e_k$. (The forcing
argument uses roots of norm~$2$ and does not apply to $Ls$ or $L\bar s$, whose
minimal vectors have norm~$4$; and indeed those two stabilizers do contain
non-monomial elements.)
(iii)~The second pass re-enumerated all three from scratch, by its own shell
and triple search, and obtained the same three orders.

\begin{remark}[A prior expectation refuted]\label{rem:12096}
It is natural to expect the order-$12096$ group $G_2(2)$ to appear as the
stabilizer of the integral-octonion order~$L$. In the paper's normalization it
does not. $L$ has minimal norm~$2$ and does not contain~$e_0$, so it is not a
unital maximal order in these coordinates; its octonion-automorphism
stabilizer has order~$1344$ and consists entirely of monomial signed
permutations. It is $L\bar s$, not~$L$, whose stabilizer is the full
$G_2(2)$. The asymmetry between $Ls$ (order $168$) and $L\bar s$ (order
$12096$) is the sharpest single fact in this computation.
\end{remark}

\subsection{The pair set $K$ and the group}

\begin{theorem}\label{thm:main}
$|K| = 6$, every element of $K$ has $A = B$, and
\[
  \Aut(\LL, +, \star) \;\cong\; C_6 \times S_3, \qquad
  \bigl|\Aut(\LL,+,\star)\bigr| \;=\; 36 .
\]
\end{theorem}

The diagonal part of~$K$ is identified by a proof: when $A = B$, $h_{A,A}$ is
just the blockwise map $(x,y,z) \mapsto (Ax,Ay,Az)$, and since $A$ is an
octonion-algebra automorphism it commutes with the sublattice conditions of
\eqref{eq:wilson} in the sense that $h_{A,A}(\LL) = \LL$ holds if and only if
$A$ stabilizes each of $L$, $L\bar s$ and $Ls$. So the diagonal part of $K$ is
exactly the triple intersection $\Stab(L) \cap \Stab(Ls) \cap \Stab(L\bar s)$,
which the enumeration above gives as a $C_6$. What is \emph{computed} is that
$K$ has no other elements: no pair with $A \neq B$ survives. The lattice kills
every genuinely asymmetric pair, even though the ambient algebra permits them
freely (Theorem~\ref{thm:ambient} has two independent $G_2$ factors).

The search over all $168 \times 12096 = 2{,}032{,}128$ candidate pairs was done
by two independent methods, which return the identical $6$ pairs
(\texttt{python\_project/src/verify\_aut\_lambda\_star.py}):

\begin{enumerate}[label=\textbf{M\arabic*.}, leftmargin=*]
\item \emph{The mod-3 glue code.} $\LL_0 := (\LL \cap D) \oplus (\LL \cap T)$
  has index $[\LL : \LL_0] = 6561 = 3^8$ in~$\LL$, and the quotient
  $\LL/\LL_0$ is an $8$-dimensional $\FF_3$-subspace of $\FF_3^{24}$. Any
  $h_{A,B}$ with $A \in \Stab(Ls)$, $B \in \Stab(L\bar s)$ preserves $\LL_0$
  automatically, so $h_{A,B}(\LL) = \LL$ if and only if the induced
  $\FF_3$-linear map $\mathrm{diag}(M_A, M_B, M_B)$ preserves the glue code.
  Matching the $168$ classes against the $12096$ classes gives exactly~$6$
  pairs.
\item \emph{Direct integrality.} With no reference to $\LL_0$, no lattice
  quotient and no $\FF_3$ linear algebra: the matrix of $h_{A,B}$ in a Hermite
  $\ZZ$-basis of~$\LL$ must be integral, and that condition, tested modulo the
  common denominator~$96$ over all $2{,}032{,}128$ pairs, returns the same~$6$.
\end{enumerate}

The second pass repeated the search in its own coordinates (hashing the
$\LL$-basis integrality condition over all $2{,}032{,}128$ pairs) and again got
exactly the same $6$ pairs, all with $A = B$, forming a $C_6$ of monomial
signed permutations fixing~$e_0$.

\subsection{Generators}

Three generators suffice. Written intrinsically:

\begin{enumerate}[label=(\arabic*), leftmargin=*]
\item \textbf{The blockwise $C_6$ generator}, of order~$6$: apply
  $A_6$ to each of the three octonion blocks, where $A_6$ is the signed
  permutation of the octonion basis fixing~$e_0$ and acting by
  \[
    e_1 \mapsto -e_5, \quad e_2 \mapsto -e_7, \quad e_3 \mapsto +e_2, \quad
    e_4 \mapsto -e_4, \quad e_5 \mapsto +e_6, \quad e_6 \mapsto +e_1, \quad
    e_7 \mapsto -e_3 .
  \]
  $A_6$ is an automorphism of $(\OO, \cdot_s)$ and stabilizes $L$, $Ls$ and
  $L\bar s$. Its underlying permutation is $(1\,5\,6)(2\,7\,3)$; its cube is
  the sign flip $e_1, e_4, e_5, e_6 \mapsto -e_1, -e_4, -e_5, -e_6$, and its
  square is a monomial automorphism of order~$3$. All six elements of $K$ are
  monomial signed permutations fixing~$e_0$.
\item \textbf{The block $3$-cycle}, of order~$3$: $(x,y,z) \mapsto (y,z,x)$.
\item \textbf{The block transposition}, of order~$2$: $(x,y,z) \mapsto (y,x,z)$.
\end{enumerate}

Generators (2) and (3) generate the $S_3$ of block permutations. Under the
Fourier normal form of \S\ref{sec:fourier} they have a clean reading: the block
$3$-cycle is the cube-root-of-unity scaling
$(\hat u_1, \hat u_2) \mapsto (\zeta^2 \hat u_1, \zeta \hat u_2)$, and the block
transposition is the Fourier conjugation $\hat u_1 \leftrightarrow \hat u_2$.
That $A_6$ is not an automorphism of the \emph{untwisted} octonion product was
checked in the second pass; the $C_6$ is a feature of $\cdot_s$.

\subsection{Structure}

The three generators, converted to $24 \times 24$ integer matrices in a Hermite
$\ZZ$-basis of~$\LL$, are exported to
\texttt{aut\_lambda\_star\_gens.g}. GAP, running its own
Schreier--Sims machinery and knowing nothing of the Python enumeration, reports
(\texttt{aut\_lambda\_star.g}):

\begin{center}
\begin{tabular}{ll}
\toprule
$|G|$ & $36$ \\
$\mathrm{StructureDescription}(G)$ & \texttt{C6 x S3} \\
abelian & no \\
center & $C_6$ \\
derived subgroup & $C_3$ \\
$-I_{24} \in G$ & no \\
conjugacy classes & $18$ \\
element orders & $1\ (\times 1)$, \; $2\ (\times 7)$, \;
                 $3\ (\times 8)$, \; $6\ (\times 20)$ \\
\bottomrule
\end{tabular}
\end{center}

The second pass rebuilt the three generators from their intrinsic description
above, in \emph{standard} $\RR^{24}$ coordinates rather than in the Hermite
basis, closed the group by hand (order~$36$), and ran GAP again on its own
matrices, obtaining the same order, the same structure description, the same
center and derived subgroup, and the same exclusion of~$-I_{24}$.

All $36$ elements were also verified one at a time, in exact arithmetic and
twice independently: each preserves~$\LL$, each satisfies
$g(u \star v) = g(u) \star g(v)$ on all $576$ basis pairs (hence, by
bilinearity, everywhere), each preserves the Leech inner product, and the $36$
are pairwise distinct linear maps.

That $-I_{24} \notin \Aut(\LL,+,\star)$ is not a surprise: negation is an
isometry but is not multiplicative for a bilinear product. It is the reason the
containment in $\Co_0$ of \S\ref{sec:co0} is strict, and it is already noted in
the main paper.

\begin{remark}[A mislabeled datum, recorded]
The first pass reported the element orders as
``$1,2,2,2,3,3,3,3,3,6(\times 9)$''. That is the multiset of orders of the $18$
\emph{conjugacy class representatives}, which is what the GAP script prints, and
not the element-order multiset of the group. The correct element-order
distribution is the one tabulated above. Both are consistent with
$C_6 \times S_3$, but the former must not be copied into the paper as an
element-order list.
\end{remark}

% ===================================================================
\section{Containment in $\Co_0$}\label{sec:co0}
% ===================================================================

\begin{theorem}\label{thm:co0}
$\Aut(\RR^{24}, \star) \subseteq O(24)$, and consequently
$\Aut(\LL, +, \star) \subseteq \Co_0$.
\end{theorem}

\begin{proof}
Let $A$ be an algebra automorphism of the real octonions. The trace
$t(x) := \tfrac18 \Tr(L_x)$ and hence the conjugation
$\bar x = t(x)\,e_0 - x$ are determined by the algebra alone, so $A$ commutes
with conjugation; and $\Nrm(x)\,e_0 = x \bar x$, so $A$ preserves the
positive-definite norm form~$\Nrm$. Hence
$\Aut(\OO, \cdot_s) \subseteq O(8)$, and it is the compact real form
of~$G_2$.

By Theorem~\ref{thm:ambient} every element of $\Aut(\RR^{24}, \star)$ is
$h_{A,B}\rho_\pi$ with $A, B \in \Aut(\OO, \cdot_s)$. By
Proposition~\ref{prop:ideals}(5), $D \perp T$ in the standard Euclidean form
of $\RR^{24}$. On $D$ the map $h_{A,B}$ acts as $(a,a,a) \mapsto (Aa,Aa,Aa)$,
which is orthogonal for the induced form ($3\Nrm(a)$); on $T$ it acts
blockwise by~$B$, also orthogonal; and $\rho_\pi$ merely permutes the three
orthogonal octonion blocks. So $h_{A,B}\rho_\pi \in O(24)$.

Finally, a $\ZZ$-linear bijection of~$\LL$ extends uniquely to $\RR^{24}$ and
its $\star$-multiplicativity extends by bilinearity, so
$\Aut(\LL,+,\star) \subseteq \Aut(\RR^{24},\star) \subseteq O(24)$. Together
with $g(\LL) = \LL$ this gives
$\Aut(\LL,+,\star) \subseteq O(24) \cap \mathrm{GL}(\LL)
 = \Aut(\LL, +, \ip{\cdot}{\cdot}) = \Co_0$.
\end{proof}

Computational confirmation: all $36$ elements were tested against the Gram
matrix of the $\LL$-basis and all preserve $\ip{u}{v}$ exactly; in the second
pass, in standard coordinates, all $36$ satisfy $g^{\!\top} g = I_{24}$
(\texttt{verify\_aut\_lambda\_star.py}, and GAP).

Note where the weight of the argument sits. Compactness of
$\Aut(\RR^{24},\star)$ is what makes $\Aut(\LL,+,\star)$ finite in the first
place (Corollary~\ref{cor:finite} gives an independent finiteness proof via the
lattice stabilizers), and it is what forces the isometry property. It is load
bearing and cannot be replaced by the trace form.

\subsection{The trace form does not do this job}

It is tempting to argue that any $\star$-automorphism preserves the
$\star$-intrinsic trace form
$B(u,v) := \Tr(L_u L_v)$, and that $B$ is (a multiple of) the Leech form.
The second half is false. Computed exactly
(\texttt{verify\_star\_algebra\_structure.py}), $B$ is diagonal with entries
$+24$ on the three $e_0$-coordinates and $-24$ on the twenty-one imaginary
coordinates:
\[
  B(u,v) \;=\; 24 \sum_{\alpha \in \ZZ/3} \Rea\bigl(u_\alpha \cdot_s v_\alpha\bigr),
  \qquad \text{signature } (3, 21).
\]
It is not a scalar multiple of the Leech Gram matrix, and it is not even of the
shape $\alpha \ip{u_D}{v_D} + \beta \ip{u_T}{v_T}$: it is blind to the $D/T$
split and separates real from imaginary parts instead. Preservation of~$B$
therefore implies neither preservation of the Leech form nor finiteness of the
automorphism group. The trace-form route fails, even though its intended
conclusion is true. This was confirmed independently in the second pass.

\subsection{Relation to what the main paper (v6) asserts}
\label{sec:v6}

The v6 Conclusion lists as an open problem: ``The automorphism group of
$(\LL, +, \star)$, equivalently the stabilizer of the product tensor
inside~$\Co_0$, and its relationship to the Conway group. Empirically,
$S_3 \subseteq \Aut(\LL, +, \star) \subsetneq \Co_0$: the block-permutation
$S_3$ acts by automorphisms, and the inclusion in $\Co_0$ is strict because
$-I_{24} \in \Co_0$ does not preserve $\star$. Full identification of
$\Aut(\LL, +, \star)$ is open.'' Against the present calculation:

\begin{itemize}[leftmargin=*]
\item \textbf{``Full identification is open'' is now closed.}
  $\Aut(\LL,+,\star) \cong C_6 \times S_3$, of order~$36$, with the three
  explicit generators of \S\ref{sec:lattice}. The open-problem list should be
  updated.
\item \textbf{$S_3 \subseteq \Aut(\LL,+,\star)$ is correct}, and the $S_3$ is
  exactly the block-permutation subgroup. It is a proper subgroup of index~$6$:
  the blockwise $C_6$ was missed.
\item \textbf{The containment $\Aut(\LL,+,\star) \subseteq \Co_0$ is correct
  and is now a theorem}, not an empirical observation. But the paper's phrasing
  presents it as if it were free: the phrase ``equivalently the stabilizer of
  the product tensor inside~$\Co_0$'' already presupposes that every
  $\star$-automorphism of $\LL$ is an isometry. That presupposition is true, by
  Theorem~\ref{thm:co0}, and the only proof we have of it goes through the
  compactness of $\Aut(\OO,\cdot_s)$. If the sentence stays, it needs that
  justification attached.
\item \textbf{The strictness $\subsetneq \Co_0$ is correct}, and the reason
  given ($-I_{24}$ is not multiplicative) is correct. It is a considerable
  understatement: $36$ against $|\Co_0| = 8{,}315{,}553{,}613{,}086{,}720{,}000$.
\end{itemize}

% ===================================================================
\section{Proved, computed, and corrected}\label{sec:corrections}
% ===================================================================

\subsection{What is proved and what is computed}

In the table, \textsc{proof} means a mathematical argument given above;
\textsc{exact} means a finite computation in exact rational or integer
arithmetic which, being exhaustive on a basis of a bilinear structure or on a
finite set, settles the statement outright; \textsc{gap} means an independent
group-theoretic identification.

\begin{center}
\small
\begin{longtable}{p{0.60\textwidth}p{0.10\textwidth}p{0.22\textwidth}}
\toprule
Statement & Status & Where \\
\midrule
\endhead
Routing identity, Lemma~\ref{lem:routing} & \textsc{proof} & \S\ref{sec:ideals}, also \textsc{exact} \\
$\Sigma$ is an algebra homomorphism; $D, T$ are two-sided ideals; $D \star T = T \star D = 0$; $\RR^{24} = D \oplus T$ & \textsc{proof} & \S\ref{sec:ideals}, also \textsc{exact} \\
$\varphi : (\OO,\cdot_s) \to (D,\star)$ is an algebra isomorphism & \textsc{proof} & \S\ref{sec:ideals} \\
$(D,\star)$ is simple & \textsc{proof} & \S\ref{sec:ideals} \\
$(T,\star)$ is simple (multiplication algebra is all of $\End(T)$, dimension $256$) & \textsc{exact} & \S\ref{sec:ideals} \\
Every $\star$-automorphism preserves $D$ and $T$; $\Aut(\RR^{24},\star) = \Aut(D,\star) \times \Aut(T,\star)$ & \textsc{proof} & \S\ref{sec:ideals} \\
Fourier identity $\widehat{(u\star v)}_k = \hat u_{2k} \cdot_s \hat v_{2k}$; $T \otimes \CC \cong$ swap-twisted algebra & \textsc{proof} & \S\ref{sec:fourier}, also \textsc{exact} over $\QQ(\zeta)$ \\
$\dim \Der(\OO,\cdot_s) = 14$, Killing form negative definite; $\dim \Der(\RR^{24},\star) = 28$ & \textsc{exact} & \S\ref{sec:ambient} \\
Component group of $\Aut(T,\star)$ is $S_3$ (centralizer of the identity component has exactly $6$ elements) & \textsc{exact} & \S\ref{sec:ambient}, Route~2 \\
$\Aut(\RR^{24},\star) = G_2 \times G_2 \times S_3$, compact, $\dim 28$ & \textsc{proof} & \S\ref{sec:ambient} \\
Every octonion-algebra automorphism is orthogonal; hence $\Aut(\RR^{24},\star) \subseteq O(24)$ and $\Aut(\LL,+,\star) \subseteq \Co_0$ & \textsc{proof} & \S\ref{sec:co0}, also \textsc{exact} on all $36$ \\
$\LL \cap T = \{(p,q,r) \in (L\bar s)^3 : p+q+r=0\}$ & \textsc{proof} & \S\ref{sec:lattice} \\
$\LL \cap D = \{(a,a,a) : a \in Ls\}$ & \textsc{exact} & \S\ref{sec:lattice} \\
$[\LL : (\LL \cap D) \oplus (\LL \cap T)] = 3^8 = 6561$ & \textsc{exact} & \S\ref{sec:lattice} \\
$A(Ls) = Ls$ and $B(L\bar s) = L\bar s$ are necessary; hence $\Aut(\LL,+,\star)$ is finite & \textsc{proof} & \S\ref{sec:lattice} \\
$\rho_\pi(\LL) = \LL$ for all six $\pi$ & \textsc{proof} & \S\ref{sec:lattice} \\
Every element of $\Stab(L)$ is a signed permutation of $e_1,\ldots,e_7$ fixing $e_0$ & \textsc{proof} & \S\ref{sec:lattice} \\
$|\Stab(L)| = 1344$, $|\Stab(Ls)| = 168$, $|\Stab(L\bar s)| = 12096$ & \textsc{exact} & \S\ref{sec:lattice} \\
Structure of the three stabilizers; $12096 \cong G_2(2) = U_3(3).2$ & \textsc{gap} & \S\ref{sec:lattice} \\
$|K| = 6$; every element of $K$ has $A = B$ & \textsc{exact} & \S\ref{sec:lattice} \\
Diagonal part of $K$ equals $\Stab(L) \cap \Stab(Ls) \cap \Stab(L\bar s)$ & \textsc{proof} & \S\ref{sec:lattice} \\
$|\Aut(\LL,+,\star)| = 36$; structure $C_6 \times S_3$; center $C_6$; derived subgroup $C_3$; $-I_{24} \notin G$ & \textsc{exact} + \textsc{gap} & \S\ref{sec:lattice} \\
Trace form $\Tr(L_u L_v) = 24 \sum_\alpha \Rea(u_\alpha \cdot_s v_\alpha)$, signature $(3,21)$, not the Leech form & \textsc{exact} & \S\ref{sec:co0} \\
\bottomrule
\end{longtable}
\end{center}

Nothing load bearing is left open. The one step that is neither an exact finite
computation nor elementary is Route~1 of \S\ref{sec:ambient}(c), the
classification of $\Aut(S)$ over~$\CC$; see Remark~\ref{rem:humanread}. It is
not load bearing, because Route~2 reaches the same conclusion by an entirely
finite and machine-checkable path.

\subsection{Corrections made during verification}

Five points where the first computational pass, or the expectation that
launched it, had to be corrected. All are recorded here rather than quietly
fixed.

\begin{enumerate}[label=\textbf{C\arabic*.}, leftmargin=*]
\item \textbf{The intrinsic-annihilator argument is circular as first stated}
  (Remark~\ref{rem:circular}). The conclusion survives; the proof is now via
  simplicity of the two ideals and a dimension count. The docstring of
  \texttt{verify\_star\_algebra\_structure.py} still carries the circular
  wording.
\item \textbf{The order-$12096$ group stabilizes $L\bar s$, not $L$}
  (Remark~\ref{rem:12096}). The prior expectation that the maximal-order
  stabilizer would be $G_2(2)$, cut down by the extra sublattice conditions, is
  wrong in the paper's normalization.
\item \textbf{The trace-form route to $\Co_0$ fails} (\S\ref{sec:co0}). Its
  conclusion is true, but for the compactness reason instead.
\item \textbf{The element-order list first reported was the list of conjugacy
  class representative orders}, not of group elements.
\item \textbf{Scale.} The $\LL$ realized by~\eqref{eq:wilson} has minimal
  norm~$8$ and Gram determinant $2^{24}$: it is $\sqrt2$ times the unimodular
  Leech lattice. This is irrelevant to $\Aut$ but no unimodularity claim should
  be attached to it.
\end{enumerate}

% ===================================================================
\section{Reproduction}\label{sec:repro}
% ===================================================================

All scripts use exact arithmetic (Python \texttt{Fraction}s or integers)
throughout; no claim in this note rests on floating-point computation. Runtimes
are on the project's reference machine. Python scripts live in
\texttt{python\_project/src/} and GAP scripts in \texttt{}; both
are named below by file name alone.

\begin{center}
\small
\begin{longtable}{>{\raggedright\arraybackslash}p{0.39\textwidth}>{\raggedright\arraybackslash}p{0.53\textwidth}}
\toprule
Script & Certifies \\
\midrule
\endhead
\texttt{verify\_ideal\_decomposition.py}
  & $\Sigma$ is an algebra homomorphism; $\RR^{24} = D \oplus T$ as mutually
    annihilating two-sided ideals; $(\OO,\cdot_s) \cong (D,\star)$; the twist
    is an algebra isomorphism onto the untwisted triple product. \\
\addlinespace
\texttt{verify\_star\_algebra\_structure.py}
  & The routing identity; $\Ann(T) = D$ and $\Ann(D) = T$; the $C_3$-Fourier
    normal form checked exactly over $\QQ(\zeta)$ on all $576$ basis pairs;
    $h_{A,B}$ and the block permutations are automorphisms;
    $\dim \Der(\OO,\cdot_s) = 14$ and $\dim \Der(\RR^{24},\star) = 28$; the
    $\star$-intrinsic trace form. About $2$ minutes. \\
\addlinespace
\texttt{verify\_aut\_lambda\_star.py}
  & The main script. $\LL \cap D$ and $\LL \cap T$; complete enumeration of the
    three octonion-automorphism stabilizers; the pair set $K$ by the two
    independent methods M1 and M2; the $36$ group elements verified one at a
    time against $\LL$, against $\star$ on all $576$ basis pairs, and against
    the Leech Gram matrix; writes the GAP generators. About $5$ minutes. \\
\addlinespace
\texttt{verify\_aut\_octonion\_crosscheck.py}
  & Independent cross-checks of the three stabilizer enumerations: re-run from a
    different generating triple, and an assumption-free brute force over all
    $645{,}120$ signed permutations for $\Stab(L)$; writes $8 \times 8$ GAP
    generators. About $9$ minutes. \\
\addlinespace
\texttt{aut\_lambda\_star.g}
  & Reads three explicit $24\times24$ integer matrices and reports
    $\mathrm{Size} = 36$, $\mathrm{StructureDescription} = $ \texttt{C6 x S3},
    center $C_6$, derived subgroup $C_3$, and $-I_{24} \notin G$. \\
\addlinespace
\texttt{octonion\_stabilisers.g}
  & Identifies $\Stab(L) = (C_2{\times}C_2{\times}C_2).\mathrm{PSL}(3,2)$ of
    order $1344$; $\Stab(Ls) = (C_2{\times}C_2{\times}C_2){:}(C_7{:}C_3)$ of
    order $168$; $\Stab(L\bar s) = \mathrm{PSU}(3,3){:}C_2$ of order $12096$,
    confirmed isomorphic to the ATLAS group $U_3(3).2 = G_2(2)$. \\
\addlinespace
\texttt{aut\_lambda\_star\_gens.g},
\texttt{octonion\_stabilisers\_gens.g}
  & Machine-generated generator files. Do not edit. \\
\bottomrule
\end{longtable}
\end{center}

The earlier probe \texttt{python\_project/src/probe\_aut\_lambda\_star.py}
recorded the $S_3$ of block permutations and the exclusion of $-I_{24}$; those
observations are subsumed here.

The independent second pass rebuilt the octonions, $\sigma$, $\cdot_s$, $L$,
$s$, $\bar s$, $Ls$, $L\bar s$ and a $\ZZ$-basis of~$\LL$ from the definitions
in \S\ref{sec:setting} (the $\LL$-basis from the defining congruences, by an
$\FF_2$-linear solve followed by Hermite reduction, rather than from a
generating set), rebuilt the three generators in standard coordinates from
their intrinsic description, and re-derived the ambient group by Route~2. It
used none of the scripts above. Its working files are not part of the
repository.

\end{document}
